\cotefile{157-2}
\scanpage{368}
\mspage{107}
% archive sheet number: 367
$C_\infty$ et $H_C$ sont ``bons'', \unsure{il|ils} nous \unsure{permettront|permettent} de décrire un \unsure{foncteur|facteur} important qui domine la situation envisagée. Moyennant une hypothèse très forte, mais \deleted{\unsure{souvent}} satisfaite dans des cas importants, on peut dans ces deux formules remplacer $\mathrm{Cat\,hot}$ par $\mathrm{Hot}$.

\deleted{Par contre, si} Par contre, sauf pour le cas où $S$ est réduit à un point et où $\mathrm{Hom}_{\mathrm{Cat\,hot}}$ ou $\mathrm{Hom}_{\mathrm{Hot}}$ de $S$ dans $C$ n'est autre que $\pi_0(C)$ $(\simeq\pi_0(C_\infty)\simeq \pi_0(H_C))$, donc effectivement ``bon'', je ne sais pas si
\[
\underset{\substack{\Vert\\\pi_0\underline{\mathrm{Hom}}(S,C)}}{\mathrm{Hom}_{\mathrm{Cat\,hot}}(S,C)}
\quad\text{ou}\quad\mathrm{Hom}_{\mathrm{Hot}}(S,C)
\]
est lié de façon raisonnable aux \deleted{\unsure{deux|deux grands|dans quelque}} ensembles précédents (déterminés par $H_C$, resp.\ \inserted{par} $C_\infty$), \deleted{et} \deleted{s'ils sont} p. ex. s'ils lui sont isomorphes de façon naturelle, et s'ils ont une signification géométrico-topologique raisonnable. Je commence à soupçonner que \uline{non} -- \unsure{alors qu'on|dans quoi on} aurait ici un cas où une relation \deleted{\unsure{entre}} de type très général concernant un $\pi_0$, ici donc
\[
\pi_0(C)\simeq\pi_0(\unsure{C|L}_\infty)\simeq\pi_0(H_C),
\]
ne se prolonge pas en un iso. sur les $\pi_i$ supérieurs, et en une \deleted{\unsure{bonne}} propriété de \unsure{hoto-équivalence|hoto-équivalences}.
\scanpage{369}
\mspage{108}
% archive sheet number: 368

Je ne trouve d'ailleurs aucune flèche directe entre $C$ et l'un des objets $C_\infty$, $H_C$ ; tout ce que je vois \deleted{\unsure{c'est}} dans ce sens, c'est
\diagram{\begin{tikzcd}[ampersand replacement=\&,column sep=large]
C_{\mathbb N^\circ}\arrow[r,"{\substack{\text{foncteur}\\\text{de locali-}\\\text{sation}}}"]\arrow[d] \& C_\infty\\
C \&
\end{tikzcd}}
où la flèche verticale est la flèche de projection de $C_{\mathbb N^\circ}=C\times_T\mathbb N^\circ$ (\unsure{avec|sur} $\mathbb N^\circ\mathrel{\unsure{\overset{\underline t}{\relbar\joinrel\relbar\mkern-2mu\mapstochar}|\overset{\underline t}{\longrightarrow}}}T$). Ceci n'établit une relation \deleted{\unsure{entre}} topologique directe entre $C$ (vu dans $\mathrm{Cat\,hot}$ ou dans $\mathrm{Hot}$ simplement) et $C_\infty$, que si les deux flèches précédentes sont \deleted{\unsure{des Cat}} tout au moins des \unsure{Hot-équivalences|Hoto-équivalences}. Et le cas-type (problématique) où la question de la relation entre $C$ et $C_\infty$ se pose, est donc celui où $T=\mathbb N^\circ$, i.e. celui d'un système inductif quelconque de \deleted{\unsure{(petites)}} \inserted{petites} catégories, indexé par $\mathbb N$, et interprété comme défini par une cat. fibrée $C=C_{\mathbb N^\circ}$ sur $\mathbb N^\circ$. Dans ce cas particulier, on a du moins une flèche canonique
\[
C\longrightarrow C_\infty
\]
sur laquelle on peut se poser la question si c'est
\scanpage{370}
\mspage{109}
% archive sheet number: 369
une $\mathrm{Hot}$-équivalence. Sinon, cela prouvera que dans le cas général, le \deleted{type} $\mathrm{Hot}$-type de $C$ (et a fortiori \deleted{\unsure{la}}\unsure{sa|la} classe d'isom. dans $\mathrm{Cat\,hot}$) n'a pas de sens géométrique intéressant. \deleted{\unsure{C}} C'est à dire que $C$ est intéressant uniquement en \uline{\unsure{tant} que catégorie fibrée sur $T$}, permettant \inserted{notamment} de construire l'\inserted{\unsure{ind-}}objet $H_T$ de $\mathrm{Cat\,hot}$, ou de $\mathrm{Hot}$, mais pas en tant que ``catégorie nue''.

Mais je viens de me convaincre que $C_{\mathbb N^\circ}\mathrel{\unsure{-|\to}}C_\infty$ est bien une $\mathrm{Hot}$-équivalence. Donc la flèche $C_{\mathbb N^\circ}\longrightarrow C$ définit une flèche \uline{dans $\mathrm{Hot}$}
\[
C_\infty\longrightarrow C
\]
(en composant avec l'inverse \inserted{dans $\mathrm{Hot}$} de la précédente). Est-ce un iso ? Donc pour cette question, on est ramené à la question pour $C_{\mathbb N^\circ}\longrightarrow C$. Mais si on prend $C=T$, cela pose déjà la même question pour $\underline{\mathbb N}^\circ\longrightarrow\deleted{\unsure{C_\infty}}T$, donc la question si $T$ est asphérique, sous la seule hypothèse que $T_0$ soit filtrante à g. et dénombrable à l'infini \inserted{(à g.)}. Mais bien sûr c'est faux (p. ex. si $T$ n'a qu'un seul objet, il n'est pas pour autant asphérique !).

Cela me convainc donc qu'à défaut
\scanpage{371}
\mspage{\deleted{\unsure{109}}\unsure{110|119}}
% archive sheet number: 370
d'hypothèses supplémentaires draconiennes sur $\deleted{\unsure{C}}T$, la catégorie $C$ ``nue'' n'offre pas d'intérêt particulier, \inserted{à part son $\pi_0$ (qui est \unsure{donné|déterminé} aussi par $C_\infty$ et par $H_{\unsure{C|T}}$).} Il n'est pas dit, \deleted{\unsure{certes}}, que $\unsure{\underline{\mathrm{Tg}}|\underline{T_g}}$ ne satisfasse pas à des ``conditions draconiennes''. Mais comme $C$ ``nu'' ne \deleted{se} s'insère pas dans un formalisme sympa (p. ex. on n'a pas d'interprétation géométrique de $\mathrm{Hom}_{\mathrm{Cat\,hot}}(S,C)$, ou de $\mathrm{Hom}_{\mathrm{Hot}}(S,C)$), je crois que pour la suite je peux en faire abstraction. Exit \unsure{$\to$|a)} !

\hfill$\backslash$\hfill\mbox{}

\hfill 5.11.90

J'ai quand même passé la journée entière hier à essayer de comprendre la relation de $C$ et de son type dans $\mathrm{Hot}$, à celui de $C_\infty$. En effet, comme on a finalement une flèche dans $\mathrm{Hot}$
\[
C_\infty\longrightarrow C,
\]
déduite du diagramme dans $(\mathrm{Hot})$
\diagram{\begin{tikzcd}[ampersand replacement=\&]
C_{\mathbb N^\circ}\arrow[r,"\sim"]\arrow[d]\& C_\infty\\
C\&
\end{tikzcd}}
de prouver que $C_\infty\longrightarrow C$ est un isom. dans $\mathrm{Hot}$
\scanpage{372}
\mspage{\deleted{\unsure{110}}\unsure{111|110}}
% archive sheet number: 371
semble la façon la plus simple, et peut-être unique, de prouver que l'image de $C_\infty$ dans $\mathrm{Hot}$ ne dépend pas, à isom. unique près, des choix faits pour construire $C_\infty$. Sans compter que dans le cas \unsure{motivant}, $C=\underline{\mathrm{Ch}}_{\unsure{\Sigma|Z}}(X)$, ou $\underline{\mathrm{Ch}}_{\Sigma}(X;a,b)$, $T=\unsure{\underline{\mathrm{Tg}}|\underline{T_g}}$, \inserted{c'est} la catégorie \unsure{des|de} chemins $\deleted{\unsure{C=\mathrm{Ch}}}$ $C$ qui donne lieu à une \deleted{\unsure{seule}} loi de composition \inserted{des chemins} impeccable, fonctorielle, associative et tout, alors qu'il n'en est rien pour $C_\infty$. On ne trouve pas \deleted{\unsure{de}} une loi de composition
\[
C_\infty\mathbin{\unsure{\times|<}}C_\infty\mathrel{\unsure{-|\to}}C_\infty
\]
dans \deleted{\unsure{Hot,}} $\mathrm{Cat}$, \deleted{\unsure{ou}} ni même dans $\mathrm{Cat\,hot}$ ou dans $\mathrm{Hot}$, ou plutôt, si on arrive \inserted{\unsure{pourtant|péniblement}} à \deleted{définir} construire \inserted{une flèche} $C_\infty\times C_\infty\longrightarrow C_\infty$ dans $\mathrm{Cat}$, c'est encore \inserted{laborieusement,} à \unsure{coup|coups} de choix $\pm$ arbitraires, \unsure{cran par cran|cas par cas}, et sans savoir si ce qu'on obtient a une \deleted{\unsure{sig}} propriété d'\unsure{unicité}, ne serait-ce que dans $\mathrm{Hot}$ ; on ne trouve une interprétation géométrique \unsure{tangible|intelligible} \inserted{de cette composition} que \deleted{\unsure{par le}} dans $\mathrm{Hot}_0^{\unsure{\wedge|1}}$ (où on a une véritable loi de groupoïde), ou dans $H_T$ en tant qu'objet de $\mathrm{Ind\,Cat\,hot}$
\scanpage{373}
\mspage{112}
% archive sheet number: 372
ou de $\mathrm{Ind\,Hot}$. (Dont la limite inductive dans $\mathrm{Hot}_0^{\unsure{\wedge|1}}$ \unsure{n'est} autre que l'élément de $\mathrm{Hot}_0^\wedge$ défini par $C_\infty$.)

Mes efforts pour prouver que \deleted{\unsure{$C$}}
\[
C_\infty\mathrel{\unsure{-|\to}}C
\]
était un iso. dans $\mathrm{Hot}$, ou ne serait-ce qu'un $1$-iso (i.e. un isom. non seulement pour les $\pi_0$, ce qui est tjrs vrai, mais aussi pour les $\pi_1$ de \deleted{\unsure{toutes}} \unsure{les|ces} composantes), ont échoué. J'ai trouvé qu'il \unsure{suffirait|suffisait} que le foncteur
\[
T\longrightarrow T_0
\]
soit ou bien lisse et coasphérique, ou bien propre et asphérique -- donc \unsure{tout est} OK si $T=T_0$ i.e. si $T$ est ordonné. Mais dans le cas $T=\unsure{\underline{\mathrm{Tg}}|\underline{T_g}}$, le foncteur $T\mathrel{-}T_0$ n'est ni propre, ni lisse (\unsure{pas|pr} $=$ $\unsure{\mathbb D|D|0}$-propre ou $\mathbb D$-lisse).

Finalement j'en suis arrivé à la conviction que dans le cas qui nous occupe $(\deleted{\unsure{C}}T=\underline{\mathrm{Tg}},\ C=\underline{\mathrm{Ch}}_{\unsure{\Sigma_{\mathscr B}|\Sigma_\beta}}(X;a,b))$ la \deleted{foncteur} \inserted{flèche} $C_\infty\longrightarrow C$ n'est en général \uline{pas} un \unsure{iso} \unsure{dans} $\mathrm{Hot}$. En effet, dans le cas $a=b$, \deleted{cela} si cela
\scanpage{374}
\mspage{113}
% archive sheet number: 373
on aurait une description du $\mathrm{Hot}$-groupe $\Omega(X,a)=\underline{\mathrm{Ch}}_{\unsure{\Sigma|Z}}(X,a)$ en termes d'une
% slanted
\margin{On fait ici $\Sigma=\unsure{\mathrm{Fl}|\mathrm{F}\ell}(X)$} catégorie $C$ munie d'une loi de composition \uline{associative} (et unitaire) \inserted{\uline{dans $\mathrm{Cat}$}}, \unsure{relevant|redonnant} la loi de groupe de $\Omega$ en passant de $(\mathrm{Cat})$ à $\mathrm{Hot}$. Prenant le Nerf de $C$, on trouverait un \unsure{monoïde|modèle} $\frac12$ simplicial qui, dans $\mathrm{Hot}$, donnerait la loi de groupe de $\Omega(X,a)$. Or je me rappelle \unsure{qu'on m'a|qu'il m'a} expliqué qu'en général, la \unsure{loi de groupe} de $\Omega(X,a)$ ne \unsure{saurait|pourrait} pas être \unsure{donnée|représentée} par une loi \uline{associative} sur un \unsure{ensemble} $\frac12$ simplicial. Ainsi, l'insuccès de mes efforts semble tenir, non à ma maladresse, mais à la nature des choses.
% slanted
\margin{C'est une \unsure{erreur}, -- ça doit être \unsure{plein|pour|faux} d'\unsure{``autres|autres} objets groupe'' de $\mathrm{Hot}$ que les groupes de la \unsure{catég.} \unsure{J'ai vu|j'ai lu} dans \unsure{Quillen} que $\Omega(X,x)$ doit bien se décrire par un \unsure{tel} groupe $\frac12$ simplicial (quand $X$ est un ens. q.c. $\frac12$ simplicial, à un sommet !). Donc \unsure{le problème|la question} si $C_\infty\simeq C$ dans $\mathrm{Hot}$ \unsure{reste} \unsure{entier|entière} !}

Je vais résumer les faits positifs auxquels je suis arrivé dans le diagramme suivant, qui a lieu dans la situation générale $C\mathrel{\unsure{-|\to}}T$ envisagée au début du présent paragraphe :
\scanpage{375}
\mspage{114}
% archive sheet number: 374

\deleted{\diagram{\begin{tikzcd}[ampersand replacement=\&,row sep=small,column sep=small]
\deleted{\unsure{C_{t'}}}\arrow[d,"{\deleted{\unsure{u^*}}}"] \& \deleted{\unsure{(t'}}\\
\deleted{C_t} \& \deleted{\unsure{t)}}\arrow[u,"{\deleted{\unsure{u}}}"]
\end{tikzcd}}}

\diagram{\begin{tikzcd}[ampersand replacement=\&,row sep=2.7em,column sep=1.8em,arrows={line width=.45pt}]
\& C_{t'}\arrow[d,"{u^*}"]\arrow[ddl,dashed]\arrow[ddr,dashed,"{\substack{\text{(pour }u:t\mathrel{\unsure{-|\to}}t'\\\text{dans }T\text{ )}}}"] \&\\
\& C_t\arrow[dl,"{\substack{\text{flèche définie}\\\text{dans }\mathrm{Ind\,Cat\,hot}\\(\mathrm{Ind}\,\deleted{\unsure{\mathrm{Cat}}}\mathrm{Hot},\\\mathrm{Ind\,Hot}_0^{\unsure{\wedge|1}})}}"']\arrow[dr,"{\substack{\text{flèche de }\mathrm{Cat}\\(\text{donc de }\mathrm{Cat\,hot},\\\text{de }\mathrm{Hot}\\\mathrm{Hot}_0^\wedge)}}"] \&\\
{\substack{\underrightarrow{\lim}\,C_t=
\tikz[baseline=(h.base)]{\node[draw,ellipse,inner sep=2pt](h){$H_T$};}\\
\text{dans }\mathrm{Ind\,Cat\,hot}\\(\text{ou dans }\mathrm{Ind\,Hot}\\
\text{ou dans }\mathrm{Ind\,Hot}_0^\wedge)}}\arrow[dr,"{\substack{\text{flèche définie}\\\text{dans }\mathrm{Ind\,Cat\,hot}\\\text{ou dans }\mathrm{Ind\,Hot}\\\text{ou dans }\mathrm{Ind\,Hot}_0^\wedge}}"']
\& {\substack{\text{diagramme}\\\text{comm. dans}\\\mathrm{Ind\,Hot}\\\mathrm{Ind\,Hot}_0^\wedge}} \& C\\
\& |[draw,ellipse,inner sep=2pt]| C_\infty\arrow[ur,"{\substack{\text{flèche définie dans }\mathrm{Hot}\\\text{donc dans }\mathrm{Hot}_0^\wedge\\\text{en général}\\\text{pas un iso }\unsure{\simeq|=}\\\text{dans }\mathrm{Hot}_0^\wedge}}"'] \&\\[-2.1em]
\& {\substack{\Vert\\ \underrightarrow{\lim}\,H_T\\\text{dans }\mathrm{Hot}_0^\wedge}} \&
\end{tikzcd}}

Les objets importants sont $H_T$ et $C_\infty$, alors que $C$ \inserted{(sur $T$)} et les $C_t$ jouent un rôle technique intermédiaire, pour construire $H_T$ et $C_\infty$. On peut considérer $C_\infty$ comme une catégorie (objet de $(\mathrm{Cat})$), mais dont la construction dépend de nombreux \unsure{arbitraires}, à l'inverse de $H_T$ dans $\mathrm{Ind\,Cat\,hot}$. \unsure{C'est} l'objet de $\mathrm{Hot}_0^\wedge$ déduit de $C_\infty$ qui a un sens canonique, et peut être décrit comme la limite inductive de $H_T$ comme objet de $\mathrm{Ind\,Hot}_0^\wedge$, i.e. comme \unsure{représentant} le $=$ foncteur sur $\mathrm{Hot}_0$ que $H_T$.
\scanpage{376}
\mspage{115}
% archive sheet number: 375

Dans le cas géométrique $T=\unsure{\underline{\mathrm{Tg}}|\underline{T_g}}$, et \unsure{d'} $C$ de la forme $\underline{\mathrm{Ch}}_{\unsure{\Sigma|\Sigma_{\cdot}}}(X;a,b)$, donnant lieu à $C_\infty(a,b)$, on \uline{n'a pas} de définition naturelle d'un \inserted{\unsure{bi-|bi}}foncteur de composition des chemins
\[
C_\infty(a,b)\times C_\infty(b,c)\longrightarrow C_\infty(a,c).
\]

On peut définir, certes, des accouplements
\[
\underline{\mathrm{Ch}}_{\Sigma,\unsure{\tau|z}}(a,b)\times
\underline{\mathrm{Ch}}_{\Sigma,\tau'}(b,c)\mathrel{\unsure{-|\to}}
\underline{\mathrm{Ch}}_{\Sigma,\tau\tau'}(a,c),
\]
\deleted{mais} ou, avec les notations relatives à la donnée de
\[
\underline{t}:\mathbb N^\circ\mathrel{-}\deleted{\unsure{C}}\underline{\mathrm{Tg}}
\]
\[
\underline{\mathrm{Ch}}_n(a,b)\times\underline{\mathrm{Ch}}_m(b,c)\mathrel{-}
\underline{\mathrm{Ch}}_{\unsure{n+m|n*m}}(a,c),
\]
mais ces accouplements \uline{ne sont pas} compatibles avec les morphismes de transition des trois syst. inductifs de catégories $(\mathrm{Ch}_n)_{n\in\mathbb N}$, ils ne le sont \uline{qu'à homotopie près}. Cela suffit cependant pour avoir des \unsure{accouplement|accouplements} parfaits
\[
H_T(a,b)\times H_{\unsure{T|t}}(b,c)\mathrel{-}H_T(a,c),
\]
avec associativité, et \inserted{(si $\Sigma=\mathrm{Fl}(X)$ tout entier)} l'existence de \uline{vrais} inverses, i.e. de \deleted{foncteurs} flèches (dans $\mathrm{Ind}\,\deleted{\unsure{\mathrm{Hot}}}\mathrm{Cat\,hot}$) de ``passage au chemin opposé''
\scanpage{377}
\mspage{116}
% archive sheet number: 376
\[
\tag*{$\unsure{(\alpha)|(2)}$}
H_T(a,b)\mathrel{\unsure{-|\to}}H_T(b,a),\quad c\mathrel{\unsure{\mapstochar\relbar|\mapsto}}c^\circ,
\]
donnant lieu au diagramme commutatif exprimant que c'est bien un inverse pour la composition. Une autre façon de le dire, c'est que pour \unsure{$\mathrm V{\Rightarrow}$|tout} objet $Y$ de $\mathrm{Cat\,hot}$ (ou de $\mathrm{Ind\,Cat\,hot}$), \deleted{\unsure{que}} les ensembles
\[
\mathrm{Hom}_{\mathrm{Cat\,hot}}(Y,H_T(a,b))\overset{\text{\unsure{def|déf}}}{=}H_T^Y(a,b)
\]
forment, pour $a,b\in X$ variables, \inserted{les flèches d'}un \uline{groupoïde}$^{*}$ \inserted{$\mathrm{Ch}_{\Sigma}^{Y}(X)$} ayant pour \inserted{ensemble de} sommets $\mathrm{Ob}\,X$, et $H_T^Y(a,b)$ étant formé des flèches de $\unsure{a_x|a}$ vers $b_{\deleted{\unsure{a}}}$, \unsure{et|les} \deleted{\unsure{fonct}} l'application
% slanted
\margin{Si $\Sigma=\mathrm{Fl}(X)$. Si $\Sigma$ est plus petit, au lieu d'un groupoïde, on trouve une \uline{catég\unsure{orie}} sans plus.}
\[
\deleted{\unsure{\mathrm{Ch}_{\Sigma}}}\ H_T^Y(a,b)\mathrel{-}H_T^Y(b,a),\qquad c\mathrel{\mapstochar\relbar}c^\circ
\]
\deleted{\unsure{fonct}} \unsure{déduite|déduit} du foncteur $(\alpha)$ plus haut n'est autre que l'application de passage à l'inverse. Cette \deleted{groupoïde} \inserted{catégorie} n'est d'ailleurs autre que celle déduite de
\[
X(Y)\Sigma_{\unsure{Y|y}}^{-1}\qquad\text{où}\quad
\left\{\begin{aligned}
&X(Y)=\underline{\mathrm{Hom}}_{\mathrm{Cat}}(Y,X)\\
&\Sigma_Y\subset\mathrm{Fl}(X_Y)\quad\text{défini}\\
&\qquad\text{en termes de }\Sigma\\
&\qquad\text{``argument par}\\
&\qquad\text{argument'',}
\end{aligned}\right.
\]
à l'aide du \deleted{\unsure{$\mathrm{fl.}$}} foncteur
\[
X\mathrel{-}X_Y,\qquad a\mathrel{\mapstochar\relbar}a_Y
\]
\scanpage{378}
\mspage{\deleted{\unsure{116}}117}
% archive sheet number: 377
associant à \deleted{\unsure{tout|\ensuremath{\mapsto}}} $a\in X$ le foncteur \uline{\unsure{const.}} \inserted{sur $Y$} de valeur $a$, \unsure{en|on} \deleted{\unsure{prenant l'image inverse du}} \inserted{\unsure{prenant|posant}} \deleted{\unsure{groupoïde}}
\[
H_T(a,b)\simeq\mathrm{Hom}_{X(Y)\Sigma_{\unsure{Y|y}}^{-1}}(a_Y,b_Y).
\]

Dans le $=$ cas géométrique, je me suis convaincu d'ailleurs que $C_\infty$ était, en tant qu'objet de $\mathrm{Hot}$ (pas seulement de $\mathrm{Hot}_0^{\unsure{\wedge|1}}$) \uline{unique à \inserted{$\mathrm{Hot}$-}isomorphisme près}, mais (chose étrange) peut-être pas à $\mathrm{Hot}$-isomorphisme \uline{unique} près.

De façon générale, revenant aux conditions du début du \S, $C_\infty$ doit s'écrire $C_\infty(\underline t)$, elle dépend du choix du foncteur cofinal
\[
\underline t:\mathbb N^\circ\longrightarrow T,
\]
Quand $\underline t$ varie dans la catégorie
\[
\underline{\mathrm{Hom\,cof}}(\mathbb N^\circ,T),
\]
les $\underline C_\infty(\underline t)$ varient de façon contravariante en $\underline t$ (donc définissent une catégorie fibrée sur $\underline{\mathrm{Hom}}\ \underline{\mathrm{cof}}(\mathbb N^\circ,T)$). De plus, le foncteur
\[
C_\infty(\underline t')\mathrel{\unsure{-|\to}}\underline C_\infty(\underline t)
\]
associé à un morphisme de systèmes projectifs
\[
\underline t\mathrel{-}\underline t'
\]
est un \uuline{iso.} dans \uline{$\mathrm{Hot}$}. Donc \deleted{les} le
\scanpage{379}
\mspage{\deleted{\unsure{117}}118}
% archive sheet number: 378
le foncteur
\[
\underline t\mathrel{\unsure{\mapstochar\relbar|\mapsto}}C_\infty(\underline t),\quad
\underline{\mathrm{Hom\,cof}}(\mathbb N^\circ,T)\longrightarrow(\mathrm{Hot})
\]
se factorise par le groupoïde fondamental
\[
\Pi\bigl(\underline{\mathrm{Hom\,cof}}(\mathbb N^\circ,T)\bigr)\mathrel{\unsure{-|\to}}(\mathrm{Hot}).
\]

Quand ce groupoïde, i.e. \inserted{quand} la catégorie $\underline{\mathrm{Hom\,cof}}$ elle-même est connexe, on peut \unsure{donc|dé-} conclure que tous les $C_\infty(\underline t)$ sont isomorphes entre eux dans $\mathrm{Hot}$. Pour conclure même à un système \uline{transitif} d'isom. entre les $C_\infty(\underline t)$, i.e. que $C_\infty(\underline t)$ est ``unique à isom. \uline{unique} près'', il faudrait savoir \inserted{de plus} que ce groupoïde fondamental est \deleted{\unsure{simplement c}} équivalent à la catégorie ponctuelle (groupe fondamental trivial), ou, comme on dit, que $\underline{\mathrm{Hom}}\ \underline{\mathrm{cof}}$ soit \uline{simplement connexe}, ou $1$-connexe.

La connexité de $\underline{\mathrm{Hom\,cof}}$ est assurée quand \inserted{la catégorie} $T$ a (comme c'est le cas pour $\unsure{\underline{\mathrm{Tg}}|\underline{T_g}}$), une loi de composition $t,t'\mathrel{\unsure{\to|-}}t*t'$, avec des morphismes fonctoriels en $t,t'$
\diagram{\begin{tikzcd}[ampersand replacement=\&,row sep=small,column sep=normal]
\&t\\[-.6em]
t*t'\arrow[ur]\arrow[dr]\&\\[-.6em]
\&t'
\end{tikzcd}}
(cette dernière condition étant \unsure{automatique}
\scanpage{380}
\mspage{119}
% archive sheet number: 379
si la loi de composition admet l'objet final de $T$, supposé exister, comme unité bilatère). Dans ce cas, $\underline t=(t_n)$ et $\underline t'=(t'_n)$ sont reliés entre eux par $\underline t*\underline t'=(t_n*t'_n)_n$, \deleted{\unsure{et}} \inserted{puisqu'}on a un morphisme de systèmes projectifs
\diagram{\begin{tikzcd}[ampersand replacement=\&,row sep=normal,column sep=small]
\&\underline t*\underline t'\arrow[dl]\arrow[dr]\&\\
\underline t\&\&\underline t'\,.
\end{tikzcd}}

Je ne sais d'ailleurs pas si la flèche dans $\mathrm{Hot}$
\[
C_\infty(\underline t)\mathrel{\unsure{-|\to}}C_\infty(\underline t')
\]
associée à $\underline t'\mathrel{-}\underline t$ dans $\underline{\mathrm{Hom\,cof}}$ \deleted{\unsure{de}} ne dépend que de $\underline t_{\unsure{0|\circ}}$, $\underline t'$, et pas du choix de $\underline t'\longrightarrow\underline t$. Ce qui est clair, \deleted{\unsure{c'est que dans le cas}} \inserted{c'est que son image} dans $\deleted{\unsure{T=\underline{\mathrm{Tg}}|T=\underline{T_g}}}$ $\mathrm{Hot}_0^{\unsure{\wedge|1}}$ est unique, c'est l'identité quand on identifie $C_\infty(\underline t)$ et $C_\infty(\underline t')$ au seul et $=$ élément de $\mathrm{Hot}_0^\wedge$ défini par $\underline{\lim}\,H_T$. \deleted{On doit avoir}

Et pourtant, je m'aperçois à l'instant que, sous l'hypothèse de l'existence de la loi de composition et de son unité bilatère, et admettant de plus que la loi en question est associative à isom. près, \uline{on a transitivité} dans les iso. (canoniques !) entre les $\underline C_\infty(t)$ qu'on vient d'indiquer !
\scanpage{381}
\mspage{120}
% archive sheet number: 380
(Sans avoir pour autant à calculer le groupe fondamental de $\underline{\mathrm{Hom}}\ \underline{\mathrm{cof}}(\mathbb N^\circ,T)$ !) Si on a trois éléments $\underline t,\underline t',\underline t''$ de $\underline{\mathrm{Hom}}\ \underline{\mathrm{cof}}$, il suffit de regarder le diagramme déduit du diagramme commutatif suivant
\diagram{\begin{tikzcd}[ampersand replacement=\&,row sep=small,column sep=small]
\&\&t*t'*t''\arrow[dl]\arrow[dr]\&\&\\
\&\underline t*\underline t'\arrow[dl]\arrow[dr]\&\&t'*t''\arrow[dl]\arrow[dr]\&\\
\underline t\&\&\underline t'\&\&\underline t''
\end{tikzcd}}
en lui appliquant le contrafoncteur $\underline s\mathrel{\unsure{\mapstochar\relbar|\mapsto}}C_\infty(\underline s)$ de $\underline{\mathrm{Hom\,cof}}$ dans $\mathrm{Hot}$.

Donc dans le cas \deleted{des} \inserted{$C(a,b)={}$}$\underline{\mathrm{Ch}}_{\unsure{\Sigma|Z}}(X;a,b)$ \deleted{\unsure{$\underline{\mathrm{Ch}}_\Sigma(X)$}}, je n'ai plus aucun scrupule à considérer la catégorie $C_\infty(a,b)$ comme un objet \uline{canonique} de $\mathrm{Hot}$, relié à \deleted{\unsure{un|au}} l'ind-objet canonique $H_T(a,b)$ \inserted{de $\mathrm{Hot}$} par un lien canonique
\[
C_\infty(a,b)\mathrel{\leftarrow\deleted{\unsure{>}}}H_T(a,b),
\tag*{$(\unsure{\alpha|a})$}
\]
qui dans $\mathrm{Hot}_0^{\unsure{\wedge|1}}$ donne un isomorphisme
\[
C_\infty(a,b)\simeq\underline{\lim}\,H_T(a,b),
\tag*{$(\unsure{\alpha'|\alpha_0|\alpha})$}
\]
i.e. la flèche $\alpha$, induit, pour $\unsure{\zeta|\xi}\in\underline{\mathrm{Ob}}\,\mathrm{Hot}$, une \deleted{\unsure{bije}}application ensembliste
\[
\mathrm{Hom}_{\mathrm{Hot}}\bigl(\zeta,C_\infty(a,b)\bigr)
\mathrel{\leftarrow\deleted{\unsure{>}}}
\underset{\substack{\Vert\\
\underrightarrow{\lim}\,\mathrm{Hom}_{\mathrm{Hot}}(\zeta,C_n(a,b))}}
{\mathrm{Hom}_{\mathrm{Ind\,Hot}}\bigl(\zeta,H_T(a,b)\bigr)}
\]
qui est \uline{bijective} si $\zeta$ est dans $\unsure{\widetilde{\mathrm{Hot}}_0|\mathrm{Hot}_0}$.
\scanpage{382}
\mspage{121}
% archive sheet number: 381

Bien entendu, les flèches
\[
C_\infty(\underline t')\mathrel{\unsure{-|\to}}C_\infty(\underline t)
\tag*{$(*)$}
\]
dans $\mathrm{Hot}$, associées aux flèches $\underline t\mathrel{-}\underline t'$ dans $\underline{\mathrm{Hom\,cof}}(\mathbb N^\circ,T)$, sont en fait des flèches dans $(\mathrm{Cat})$ elle-même -- et $=$ des foncteurs pleinement fidèles, dans le cas ``géométrique'' et pour les deux flèches canoniques
\diagram{\begin{tikzcd}[ampersand replacement=\&,row sep=small,column sep=small]
\&\underline t*\underline t'\arrow[dl]\arrow[dr,dash]\&\\
\underline t\&\&\underline t'\,.
\end{tikzcd}}

Elles sont isomorphiques \uline{dans $\mathrm{Hot}$}, parce qu'elles \deleted{\unsure{limites inductives}} définissent un isomorphisme sur les ind-objets associés de $\mathrm{Cat\,hot}$, donc aussi pour les ind-objets associés de $\mathrm{Hot}$ et \deleted{de $\mathrm{Hot}$} dans $\mathrm{Hot}_0^{\unsure{\wedge|1}}$, d'où résulte que
\[
\mathrm{Hom}_{\mathrm{Hot}}(Y,C_\infty(\underline t'))\mathrel{-}
\mathrm{Hom}_{\mathrm{Hot}}(Y,C_\infty(\underline t))
\]
associé à $(*)$ est une \deleted{\unsure{isom.}} \inserted{bij.}, \uline{si $Y$ est finie}$^{*}$, cette flèche s'identifiant alors à
% slanted
\margin{Ce qui suffit à assurer que $(*)$ est iso dans $\mathrm{Hot}$\unsure{, ;| ;} $\mathrm{Hot}\to\mathrm{Hot}_0^\wedge$ est \uline{conservatif} !}
\[
\underset{n}{\underrightarrow{\lim}}\,\mathrm{Hom}_{\mathrm{Hot}}(Y,C_{t'_n}\deleted{\unsure{(t'_n)}})
\longrightarrow\underset{n}{\underrightarrow{\lim}}\,\mathrm{Hom}_{\mathrm{Hot}}(Y,C_{t_n}\deleted{\unsure{(t_n)}}).
\]
(Car \inserted{$Z\mapsto$}$\mathrm{Hom}_{\mathrm{Hot}}(Y,Z)$ commute aux limites inductives filtrantes, si $Y$ est une catégorie \uline{finie}. \deleted{\unsure{)}} Cela provient du fait qu'une $\underrightarrow{\lim}$ filtrante de catégories de \unsure{Thomason-Kan|Thomason Kan} est itou\ldots)

Mais j'ignore par contre si $(*)$ est iso dans $\mathrm{Cat\,hot}$. \unsure{Sachant|Seulement} que pour une catégorie \uline{finie} $Y$, elle induit
\[
\mathrm{Hom}_{\mathrm{Cat\,hot}}(Y,C_\infty(\underline t'))\simeq
\mathrm{Hom}_{\mathrm{Cat\,hot}}(Y,C_\infty(\underline t)),
\]
les deux membres s'identifiant alors à
